% Vietnamese translation for OI Public Library
% Source-File: IOI中国国家候选队论文/2004/金恺.ppt
\documentclass[11pt]{article}
\usepackage{oipl-vn}

\newcommand{\Oh}{\mathcal{O}}

\title{Phương pháp giới hạn trong tối ưu hình học\\{\large Ghi chú theo slide}}
\author{Jin Kai}
\date{IOI 2004}

\begin{document}
\maketitle

\origtitle{The limit method in geometric optimization}
\sourcepath{IOI 2004 / Jin Kai.ppt}

\section{Ý tưởng}

Phương pháp giới hạn dùng một thay đổi vô cùng nhỏ để chứng minh rằng nhiều
cấu hình không thể tối ưu. Nếu một điểm, đường thẳng hoặc cạnh có thể dịch
chuyển rất nhỏ mà làm hàm mục tiêu tốt hơn, cấu hình hiện tại bị loại. Nhờ đó
ta thường giảm bài toán liên tục về một tập hữu hạn ứng viên.

\section{Các mốc trên slide}

Slide nhắc các điểm \(V_1(x_1,y_1),V_2(x_2,y_2),\ldots,V_n(x_n,y_n)\) và
đường thẳng
\[
  l:\ ax+by+c=0,\qquad a^2+b^2\ne 0.
\]
Mục tiêu có dạng \(\min f(l)\). Khi thay \(l\) bằng một đường \(l'\) hoặc
\(l''\) rất gần, hiệu \(f(l')-f(l)\) cho biết hướng cải thiện. Nếu hướng cải
thiện tồn tại, đường hiện tại không phải nghiệm tối ưu.

Một mốc khác là độ phức tạp \(\Oh(n^3)\), thường xuất hiện sau khi đã chứng
minh rằng nghiệm tối ưu phải đi qua một số điểm hoặc cạnh đặc biệt, rồi liệt
kê các ứng viên hữu hạn.

\section{Cách áp dụng}

Trong bài tối ưu hình học, hãy tìm các điều kiện biên bắt buộc: cạnh đi qua
điểm nào, đường tựa chạm đa giác ở đâu, hoặc góc nào không thể xoay thêm.
Chỉ sau khi có điều kiện này, việc duyệt ứng viên mới có cơ sở và không phụ
thuộc vào trực giác hình vẽ.

\end{document}
